\subsection{Parallel Vs. Serial Performance}
It can be seen that generally the parallelisation was implemented correctly as the total time during execution was reduced as the number of ``OMP'' threads increasedbelow in \fig\ref{fig:large:OMPVTOTAL}.
Interestingly, the OpenMP solver can be seen to actually become siginificantly worse after roughly 8 threads. This was hypothesised due to the the overhead between the threads when encountering for loops and importantly caching issues.
Expanding on this further, when the number of threads increased the amount of cache used did also; therefore after a certain point the executable no onger runs as efficiently as possible.
\begin{figure}[H]
        \centering
        \includegraphics[width=0.45\textwidth]{Images/\largedomain___OMPVTOTAL.png}
        \caption{Plot showing performance of OMP implementation against serial case}
        \label{fig:large:OMPVTOTAL}
\end{figure}
As before, the normalised time plot has been produced below in \fig\ref{fig:both:OMPVNORMAL}.
\begin{figure}[H]
    \centering
    \begin{subfigure}[b]{0.45\textwidth}
        \centering
        \includegraphics[width=\textwidth]{Images/\largedomain___OMPVNORMAL.png}
        \caption{OMP (Large)}
        \label{fig:large:OMPVNORMAL}
    \end{subfigure}
    \hfill
    \begin{subfigure}[b]{0.45\textwidth}
        \centering
        \includegraphics[width=\textwidth]{Images/\smalldomain___OMPVNORMAL.png}
        \caption{OMP (Small)}
        \label{fig:small:OMPVNORMAL}
    \end{subfigure}
    \caption{Plots showing normalized time scaling plot to see the effect Open MP parallelisation has on both ``Large'' and ``Small'' domains}
    \label{fig:both:OMPVNORMAL}
\end{figure}
\vspace{-0.75cm}
It is apparant that when smaller domain settings were used the solver did not show very good scaling. This is incontrast to the MPI implementation where, even at ``Small'' settings the running time still reduced.
Specifically, there was a 48\% reduction in ``Solve'' time when increasing the serial case to have up to 8 threads. Furthermore, unlike the MPI case, the ``Write''  time was severely increased with thread number showing major overhead in communication with threads hanging about and not doing work.
This makes sense however as there was a MPI barrier within this function that would force sequential writes for the ordered output.
\subsection{Design Methodology}
When implementing Open MP for the initial design, one major design concept was used.
THis utilised the ``\#pragma omp parallel for schedule(static)'' pattern on almost all for loops where possible, ensuring no race conditions or false sharing occured.
This process however clearly was not the most effective as the performance did not scale as well as the MPI case with thread count ( and rank count).
